Koperasi di Indonesia menghadapi permasalahan struktural berupa tata kelola yang opak, pencatatan keuangan manual yang rawan manipulasi, partisipasi anggota yang terbatas, dan pengawasan yang lemah akibat ketiadaan akses \textit{real-time} terhadap data operasional. Di sisi lain, teknologi blockchain dan \textit{Decentralized Autonomous Organization} (DAO) menawarkan transparansi dan otomatisasi melalui \textit{smart contract}, namun model DAO yang ada menggunakan \textit{token-weighted voting} yang bertentangan dengan prinsip satu-anggota-satu-suara sebagaimana diamanatkan dalam Undang-Undang Nomor 25 Tahun 1992 tentang Perkoperasian. Penelitian ini bertujuan untuk merancang, mengimplementasikan, dan mengevaluasi sistem koperasi berbasis blockchain---Lakomi---yang memetakan 26 ketentuan UU 25/1992 ke dalam \textit{smart contract} yang dapat dieksekusi pada jaringan DChain.

Penelitian mengadopsi metodologi \textit{Design Science Research} (DSR) dengan tiga siklus: relevansi, rigor, dan desain. \textit{Artifact} yang dihasilkan berupa empat \textit{smart contract} yang saling terintegrasi: LakomiToken (keanggotaan berbasis ERC-20 dengan \textit{non-transferability}), LakomiVault (pengelolaan tiga jenis simpanan dan distribusi SHU enam kategori), LakomiGovern (tata kelola demokratis dengan mekanisme satu-anggota-satu-suara, kuorum 67\%, pemilihan pengurus, dan veto pengawas), dan LakomiLoans (pinjaman mikro dengan \textit{loan-to-value} berbasis kontribusi dan suku bunga tetap 5\% APY). Sistem menerapkan \textit{Role-Based Access Control} dengan delapan peran yang didistribusikan ke akun terpisah untuk memastikan pemisahan kekuasaan. \textit{Frontend} dibangun menggunakan React 18 dengan TypeScript dan wagmi v2 untuk interaksi dengan \textit{smart contract}.

Pengujian dilakukan melalui 26 skenario fungsional yang memverifikasi setiap ketentuan UU 25/1992 pada lingkungan DChain. Hasil pengujian menunjukkan tingkat kepatuhan 100\%---seluruh 26 ketentuan berhasil diimplementasikan dan berfungsi sesuai spesifikasi. Analisis gas menunjukkan bahwa operasi utama koperasi memiliki konsumsi komputasi yang wajar untuk jaringan blockchain konsorsium. Perbandingan dengan koperasi konvensional dan DAO berbasis token menunjukkan bahwa Lakomi unggul dalam transparansi (\textit{distributed ledger} yang dapat diaudit publik), demokrasi (satu-anggota-satu-suara yang dijamin \textit{smart contract}), dan otomatisasi (distribusi SHU on-chain). Arsitektur \textit{dual-track} yang memisahkan hak suara dari insentif kontribusi merupakan kontribusi desain utama yang membedakan Lakomi dari sistem yang ada dan menyelaraskan prinsip demokrasi koperasi dengan keberlanjutan finansial.

\vspace{1cm}
\noindent\textbf{Kata kunci} : blockchain, koperasi, smart contract, UU 25/1992, Design Science Research
